% TOP_PROBLEM: {"id":30007520,"problem_number":"LOCAL-30007520","title":"Recession and tail-risk forecasting","category_id":21,"category":{"id":21,"name":"economics","display_name":"Economics","description":"Open economic theory statements, published research agendas, and proposed scoped specifications; question provenance and historical priority are qualified per record.","slug":"economics","order_index":21,"created_at":"2026-10-08T00:00:00Z"},"status":"research_question","external_url":null,"legacy_ids":[],"published":false,"statement":"Design calibrated recession-tail forecasts with bounded excess prediction loss when the data process has explicitly restricted structural changes.","economics":{"original_id":9,"field":"Business cycles and macro dynamics","kind":"design","formalization_basis":"adapted_research_question","source_status":"research_agenda","priority_status":"not_established","priority_note":"Original priority is not established. No dated primary statement of this precise problem was verified.","earliest_verified_reference":{"authors":["European Central Bank"],"title":"Research agenda: Forecasting, macroeconomic and financial analysis","year":null,"url":"https://www.ecb.europa.eu/pub/economic-research/research_agenda/forecasting/html/index.en.html","doi":null,"locator":"“Forecasting and related tools,” selected areas of focus","evidence_summary":"The institution explicitly lists real-time nowcasting, forecasting, forecast combinations, and density forecasts as research priorities. It does not publish a formal minimax forecasting conjecture."},"current_reference":{"authors":["European Central Bank"],"title":"Research agenda: Forecasting, macroeconomic and financial analysis","year":null,"url":"https://www.ecb.europa.eu/pub/economic-research/research_agenda/forecasting/html/index.en.html","doi":null,"locator":"“Forecasting and related tools,” selected areas of focus","evidence_summary":"The institution explicitly lists real-time nowcasting, forecasting, forecast combinations, and density forecasts as research priorities. It does not publish a formal minimax forecasting conjecture."},"formalization_note":"The ECB lists forecasting and density-forecast priorities; this bounded-regime tail-prediction problem is an adapted specification. The agenda page is undated."},"statusQualification":"Published question or research agenda verified at the cited locator. The mathematical specification is adapted for this catalogue and includes additional assumptions; source publication does not establish first-ever priority or certify this exact model remains unsolved. The ECB lists forecasting and density-forecast priorities; this bounded-regime tail-prediction problem is an adapted specification. The agenda page is undated.","statusReviewedAt":"2026-10-08"}
% FIRST_POSTED: 2026-10-08
% ENTRY_KIND: statement_only
% !TeX program = lualatex
\documentclass[11pt]{article}
\usepackage{fontspec}
\usepackage[a4paper,margin=25mm]{geometry}
\usepackage{amsmath,amssymb,mathtools}
\usepackage{xurl}
\usepackage[hidelinks]{hyperref}
\newcommand{\sref}[2]{\hyperlink{source:#1:#2}{[S#2]}}
\setlength{\emergencystretch}{3em}
\begin{document}
% problemId:30007520
\section*{Recession and tail-risk forecasting}\label{30007520}
\subsection{Definitions and mathematical statement}
Let \(X_t\) be a vector of real-time information available at date \(t\), including data vintages rather than later revisions. Define a recession-tail event \(D_{t,H}=1\{Y_{t+H}/Y_t-1<-q\}\), where \(Y_t\) is real output, \(H\) a fixed horizon, and \(q>0\) a predetermined threshold. A forecaster produces \(p_t=f_t(X_{\le t})\in[0,1]\). Specify a family \(\mathcal P\) of data-generating processes allowing at most \(K\) regime changes, bounded parameter drift between changes, and a minimum event frequency; these restrictions must be stated numerically before evaluation.
Construct a forecast rule whose excess expected proper-score loss is small relative to an oracle that knows the current regime. For Brier loss the target is a bound on
\[
\sup_{P\in\mathcal P}E_P\sum_{t=1}^T
\big[(p_t-D_{t,H})^2-(p_t^*-D_{t,H})^2\big],
\]
where \(p_t^*=P(D_{t,H}=1\mid X_{\le t},\text{current regime})\).
Evaluate simultaneous calibration and useful discrimination for rare downturns using strictly sequential holdouts spanning structural changes. Account for overlapping horizons and revised recession labels. Establish which information and stability restrictions make improvement possible; no uniformly accurate prediction is demanded for arbitrary unforeseeable regime changes. Report uncertainty rather than selecting a model on the same crises used to assess performance.

\par\textbf{Formulation and scope.} The ECB lists forecasting and density-forecast priorities; this bounded-regime tail-prediction problem is an adapted specification. The agenda page is undated.

\par\textbf{Open-question provenance.} An explicit question or research agenda was verified at the source locator. This does not by itself certify that every later version remains unresolved \sref{30007520}{1}.

\par\textbf{Historical priority.} Original priority is not established. No dated primary statement of this precise problem was verified.

\subsection{Short English statement}
Design calibrated recession-tail forecasts with bounded excess prediction loss when the data process has explicitly restricted structural changes.

\subsection{Sources}
\begin{itemize}\raggedright
\item[\textnormal{[S1]}]\hypertarget{source:30007520:1}{} European Central Bank. Year not verified. Research agenda: Forecasting, macroeconomic and financial analysis. “Forecasting and related tools,” selected areas of focus.
\url{https://www.ecb.europa.eu/pub/economic-research/research_agenda/forecasting/html/index.en.html}.
{\small\textit{Earliest verified question/agenda reference; historical priority qualified below. The institution explicitly lists real-time nowcasting, forecasting, forecast combinations, and density forecasts as research priorities. It does not publish a formal minimax forecasting conjecture.}}
\end{itemize}
\end{document}
